% !PRINTED 709
\weberpage{709}
%<<CONT
von denen hiernach die Bestimmung der Classenzahl noch abhängt, lassen
sich durch endliche Ausdrücke darstellen. Es ist nämlich
\[
\frac{1}{n} = \int_0^1 x^{n-1}\,dx,
\]
und demnach
\[
X = \int_0^1 \Sigma \chi(n)\,x^{n-1}\,dx.
\]

Nun bleibt $\chi(n)$ ungeändert, wenn $n$ um ein Vielfaches von $m$ wächst.
Versteht man aber unter $t$ den kleinsten positiven Rest von $n$ und setzt
\[
n = t + ml,
\]
so ist $\chi(n) = \chi(t)$, und es folgt:
\[
X = \int_0^1 \overset{t}{\Sigma} \chi(t)\,x^{t-1} \underset{0,\,\infty}{\overset{l}{\Sigma}} x^{ml}\,dx.
\]

Setzt man jetzt
\weq{2}{f(x) = \overset{t}{\Sigma} \chi(t)\,x^t,}
so ist $f(x)$ eine ganze Function von $x$ vom Grade $m-1$, und zugleich ist
wegen der Relation $\Sigma \chi(t) = 0$:
\weq{3}{f(0) = 0,\quad f(1) = 0,}
also $f(x)$ durch $x(1-x)$ theilbar. Für $X$ ergiebt sich dann nach der
Summenformel für die geometrische Reihe $\Sigma x^{ml}$:
\weq{4}{X = \int_0^1 \frac{f(x)\,dx}{x(1-x^m)}.}

Um ein solches Integral zu finden, schreibt die Integralrechnung vor, den
rationalen Bruch unter dem Integralzeichen in Partialbrüche zu zerlegen.
Diese Partialbruch-Zerlegung ergiebt aber, wegen (3), wenn
\[
r = e^{\frac{2\pi i}{m}}
\]
gesetzt wird:
\[
\frac{f(x)}{x(1-x^m)} = -\frac{1}{m} \sum\limits_{1,\,m-1}^{s} \frac{f(r^s)}{x-r^s}
\]
[Bd. I, \S.\,29, (11)], und nun haben wir weiter nach bekannten
%CONT>>
